\documentclass[12pt, a4paper]{article}

% === PAQUETES ===
\usepackage[spanish]{babel}
\usepackage[utf8]{inputenc}
\usepackage[T1]{fontenc}
\usepackage{amsmath, amssymb, amsthm}
\usepackage{physics}
\usepackage[margin=2.5cm]{geometry}

\usepackage{booktabs}
\usepackage[hidelinks]{hyperref}

% === ENTORNOS ===
\theoremstyle{definition}
\newtheorem{definicion}{Definición}[section]
\newtheorem{teorema}{Teorema}[section]
\newtheorem{ejemplo}{Ejemplo}[section]
\newtheorem{observacion}{Observación}[section]

% === ENCABEZADO ===
\title{\textbf{Dioptrios: Óptica Geométrica de Superficies Refractantes}}
\author{}
\date{}

\begin{document}
\maketitle
\thispagestyle{empty}
\tableofcontents
\newpage

% ============================================================
\section{Sistemas Ópticos}
% ============================================================

\begin{definicion}[Sistema óptico]
En \textbf{óptica geométrica} se denomina \textit{sistema óptico} a un conjunto de
superficies que separan medios con \textbf{distintos índices de refracción}. Estas
superficies pueden ser \textbf{refractantes} o \textbf{espejos}.
\end{definicion}

Con frecuencia nos encontramos con sistemas formados por \textbf{superficies esféricas},
con sus centros de curvatura situados sobre una misma recta llamada \textbf{eje del
sistema} o \textbf{eje óptico}. A estos sistemas se les denomina \textit{sistemas ópticos
centrados}.

\subsection*{Clasificación}

Los sistemas ópticos se clasifican según el tipo de superficies que los forman:

\begin{itemize}
  \item \textbf{Dióptricos:} formados \textbf{solo} por superficies refractantes.
  \item \textbf{Catóptricos:} formados \textbf{solo} por espejos.
  \item \textbf{Catadiáptricos:} formados por \textbf{ambos tipos} (refractantes y espejos).
\end{itemize}

\textit{\textbf{Síntesis.}} La clasificación tripartita de los sistemas ópticos refleja la
dualidad física fundamental entre reflexión y refracción. La mayoría de los instrumentos
ópticos modernos —telescopios catadióptricos como el Schmidt-Cassegrain, por ejemplo—
explotan ambos fenómenos simultáneamente para minimizar aberraciones y reducir la longitud
del tubo óptico. El estudio del dioptrio, forma más elemental de sistema dióptrico,
constituye la base analítica de todos estos desarrollos.

% ============================================================
\section{El Dioptrio}
% ============================================================

\begin{definicion}[Dioptrio]
Un \textbf{dioptrio} es el sistema óptico formado por \textbf{una sola superficie} que
separa dos medios transparentes con \textbf{distinto índice de refracción}. Cualquier
punto del dioptrio puede tomarse como eje óptico del sistema.
\end{definicion}

\subsection{Tipos de dioptrio}

\begin{itemize}
  \item \textbf{Dioptrio esférico:} superficie esférica que separa dos medios de diferente
        índice de refracción. Según el signo del radio de curvatura: \textbf{convexo}
        (\(r > 0\)) y \textbf{cóncavo} (\(r < 0\)).
  \item \textbf{Dioptrio plano:} superficie plana que separa dos medios transparentes de
        distinto índice de refracción.
\end{itemize}

\begin{observacion}
El estudio del dioptrio esférico es fundamental porque en los espejos y en las lentes, los
componentes de los instrumentos ópticos, la luz se comporta siguiendo leyes similares.
\end{observacion}

\subsection{Notación utilizada}

A lo largo del desarrollo se emplea la siguiente notación estándar:

\begin{itemize}
  \item \(n_1\): índice del medio donde está el \textbf{objeto}
  \item \(n_2\): índice del medio donde se forma la \textbf{imagen}
  \item \(p\): distancia del objeto al vértice
  \item \(q\): distancia de la imagen al vértice
  \item \(r\): radio de curvatura de la superficie
  \item \(\theta_1,\, \theta_2\): ángulos de incidencia y refracción
\end{itemize}

\textit{\textbf{Síntesis.}} La notación adoptada es la estándar anglosajona —heredada de
la convención de Gauss— en la que el objeto, la imagen y el centro de curvatura se
referencian todos respecto al vértice óptico \(O\). Esta elección minimiza la proliferación
de signos negativos en los casos más frecuentes y unifica el tratamiento de dioptrios
esféricos y planos bajo una única ecuación.

% ============================================================
\section{Refracción en una Superficie Esférica}
% ============================================================

Considérense dos medios transparentes con índices de refracción \(n_1\) y \(n_2\), donde
la frontera entre los medios es una \textbf{superficie esférica de radio} \(r\).

Se supondrá que el \textbf{objeto} se encuentra en el medio de índice \(n_1\), en el punto
\(P\). Además, los rayos paraxiales que salen de \(P\) forman un ángulo pequeño con el eje
y entre ellos. Todos los rayos que parten del punto objeto se refractan en la superficie
esférica y los rayos refractados (o sus prolongaciones) se intersecan en un solo punto
\(Q\), que es el \textbf{punto imagen}.

\begin{definicion}[Tipos de imagen]
\begin{itemize}
  \item La imagen es una \textbf{imagen real} si está formada por la \textbf{intersección}
        de los rayos refractados.
  \item La imagen es una \textbf{imagen virtual} si está formada por la
        \textbf{prolongación} de los rayos refractados.
\end{itemize}
\end{definicion}

\begin{observacion}
Al contrario de lo que sucede en los espejos esféricos, en los dioptrios esféricos las
imágenes reales están detrás de la superficie esférica, y las imágenes virtuales se
encuentran en frente.
\end{observacion}

\subsection*{Imagen real vs.\ imagen virtual}

La siguiente tabla resume la diferencia entre espejos esféricos y dioptrios esféricos:

\begin{center}
\begin{tabular}{@{}llll@{}}
\toprule
\textbf{Característica} & \textbf{Espejo Esférico} & \textbf{Dioptrio Esférico} \\
\midrule
Imagen real    & Delante del espejo  & \textbf{Detrás} de la superficie  \\
Imagen virtual & Detrás del espejo   & \textbf{Delante} de la superficie \\
Fenómeno       & Reflexión           & Refracción                        \\
\bottomrule
\end{tabular}
\end{center}

\textbf{Elementos del dioptrio esférico:}
\(O\) = vértice óptico; \(C\) = centro de curvatura; \(p\) = distancia objeto;
\(q\) = distancia imagen; \(r\) = radio de curvatura.

\textit{\textbf{Síntesis.}} La inversión de la posición de las imágenes reales y virtuales
respecto a los espejos no es un mero detalle de convenio: refleja el hecho físico de que
la luz \textit{atraviesa} la superficie en lugar de reflejarse en ella. Esta asimetría
obliga a reinterpretar cuidadosamente el signo de \(q\) al pasar de un sistema catóptrico
a uno dióptrico.

% ============================================================
\section{Demostración de la Ecuación Fundamental}
% ============================================================

\subsection{Caso: Superficie convexa con \texorpdfstring{\(n_1 < n_2\)}{n1 < n2}
            --- imagen real}

Consideremos el caso en que \(n_1 < n_2\), utilizando la construcción geométrica que
muestra un solo rayo \(\overrightarrow{PA}\) que al refractarse en la superficie esférica
convexa de radio \(r\) intercepta al eje óptico en el punto \(Q\), que será la imagen real
de \(P\).

En la figura se identifican los siguientes elementos geométricos:
\begin{itemize}
  \item \(\theta_1\): ángulo de incidencia (entre el rayo \(\overrightarrow{PA}\) y la
        normal \(N\) en el punto \(A\))
  \item \(\theta_2\): ángulo de refracción (entre el rayo refractado y la normal)
  \item \(\alpha_1\): ángulo que subtiende el segmento \(PA\) respecto al eje óptico
  \item \(\alpha_2\): ángulo que subtiende el segmento \(AQ\) respecto al eje óptico
  \item \(\beta\): ángulo en el centro de curvatura \(C\)
  \item \(h\): altura perpendicular del punto \(A\) al eje óptico
\end{itemize}

\subsubsection*{Paso 1: Ley de Snell}

Al aplicar la \textbf{ley de Snell} a los rayos incidente y refractado en el punto \(A\):

\begin{equation}
  n_1 \sin\theta_1 = n_2 \sin\theta_2
  \label{eq:snell}
\end{equation}

\subsubsection*{Paso 2: Aproximación paraxial}

Como los ángulos \(\theta_1\) y \(\theta_2\) son pequeños (rayos paraxiales), se puede
usar la aproximación de \textbf{ángulos pequeños} (en radianes):
\[
  \sin\theta_1 \approx \theta_1 \qquad;\qquad \sin\theta_2 \approx \theta_2
\]

Con estas aproximaciones, la ley de Snell se simplifica a:

\begin{equation}
  n_1 \theta_1 = n_2 \theta_2
  \label{eq:snell_paraxial}
\end{equation}

\begin{observacion}
La aproximación paraxial es válida para rayos que viajan cerca del eje óptico. Esto
simplifica enormemente el análisis, convirtiendo las funciones trigonométricas en funciones
lineales.
\end{observacion}

\subsubsection*{Paso 3: Relaciones geométricas de los ángulos}

Se utiliza la propiedad de que un \textbf{ángulo exterior} de cualquier triángulo es igual
a la suma de los ángulos interiores no adyacentes.

Al aplicar esta propiedad a los triángulos \(\triangle APC\) y \(\triangle CAQ\) se
obtiene:

\textit{En el triángulo \(\triangle APC\):} el ángulo \(\theta_1\) es exterior al
triángulo (formado por la normal y el rayo incidente). Los ángulos interiores no
adyacentes son \(\alpha_1\) y \(\beta\):

\begin{equation}
  \theta_1 = \alpha_1 + \beta
  \label{eq:geo1}
\end{equation}

\textit{En el triángulo \(\triangle CAQ\):} el ángulo \(\beta\) es exterior al triángulo.
Los ángulos interiores no adyacentes son \(\theta_2\) y \(\alpha_2\):

\begin{equation}
  \beta = \theta_2 + \alpha_2 \qquad \Longrightarrow \qquad \theta_2 = \beta - \alpha_2
  \label{eq:geo2}
\end{equation}

\subsubsection*{Paso 4: Combinación de ecuaciones}

Sustituyendo las ecuaciones \eqref{eq:geo1} y \eqref{eq:geo2} en la ley de Snell paraxial
\eqref{eq:snell_paraxial}: \(n_1\theta_1 = n_2\theta_2\):

\begin{align}
  n_1(\alpha_1 + \beta) &= n_2(\beta - \alpha_2) \notag \\
  n_1\alpha_1 + n_1\beta &= n_2\beta - n_2\alpha_2 \notag \\
  n_1\alpha_1 + n_2\alpha_2 &= n_2\beta - n_1\beta \notag \\
  n_1\alpha_1 + n_2\alpha_2 &= (n_2 - n_1)\beta
  \label{eq:clave}
\end{align}

\begin{observacion}
La ecuación \eqref{eq:clave} es el resultado clave de la geometría: relaciona los ángulos
\(\alpha_1\), \(\alpha_2\) y \(\beta\) con los índices de refracción. El siguiente paso
será expresar estos ángulos en función de las distancias \(p\), \(q\) y \(r\).
\end{observacion}

\subsubsection*{Paso 5: Expresiones geométricas de los ángulos}

En la aproximación de ángulos pequeños: \(\tan\theta \approx \theta\). De la geometría,
los ángulos \(\alpha_1\), \(\alpha_2\) y \(\beta\) pueden expresarse en función de la
altura \(h\) del punto de incidencia \(A\) y las distancias al eje óptico:

\begin{equation}
  \alpha_1 \approx \tan\alpha_1 = \frac{h}{p}
  \label{eq:a1}
\end{equation}
\begin{equation}
  \alpha_2 \approx \tan\alpha_2 = \frac{h}{q}
  \label{eq:a2}
\end{equation}
\begin{equation}
  \beta \approx \tan\beta = \frac{h}{r}
  \label{eq:beta}
\end{equation}

\subsubsection*{Paso 6: Sustitución final}

Sustituyendo \eqref{eq:a1}, \eqref{eq:a2} y \eqref{eq:beta} en la ecuación
\eqref{eq:clave}: \(n_1\alpha_1 + n_2\alpha_2 = (n_2 - n_1)\beta\), se obtiene:

\[
  n_1 \cdot \frac{h}{p} + n_2 \cdot \frac{h}{q} = (n_2 - n_1) \cdot \frac{h}{r}
\]

Factorizando \(h\) y dividiendo ambos miembros entre \(h\) (que es distinta de cero, pues
\(A\) no está sobre el eje):

\begin{equation}
  \boxed{\frac{n_1}{p} + \frac{n_2}{q} = \frac{n_2 - n_1}{r}}
  \label{eq:fundamental}
\end{equation}

Esta es la \textbf{Ecuación Fundamental del Dioptrio Esférico} (Ec.\ 4-30).

% ============================================================
\section{Casos Particulares de Demostración}
% ============================================================

\subsection{Caso: Superficie convexa con \texorpdfstring{\(n_1 > n_2\)}{n1 > n2}
            --- imagen virtual}

Cuando \(n_1 > n_2\), la luz pasa de un medio \textbf{más denso} a uno \textbf{menos
denso}. Por la ley de Snell, el rayo refractado se aleja de la normal
(\(\theta_2 > \theta_1\)). El rayo refractado no cruza el eje, pero su prolongación hacia
atrás sí lo hace, formando una \textbf{imagen virtual}.

La ley de Snell paraxial se mantiene:
\[
  n_1 \cdot \theta_1 = n_2 \cdot \theta_2
\]

De los triángulos \(\triangle APC\) y \(\triangle CAQ\) se obtiene:

\[
  \theta_1 = \alpha_1 + \beta \qquad\text{(igual que antes)}
\]
\[
  \theta_2 = \alpha_2 + \beta \qquad\text{(ahora \(\theta_2\) es ángulo exterior en
  \(\triangle CAQ\))}
\]

\begin{observacion}
La diferencia clave: en el caso \(n_1 < n_2\) teníamos \(\beta = \theta_2 + \alpha_2\),
pero ahora \(\theta_2 = \alpha_2 + \beta\) porque la posición de \(Q\) respecto a \(C\)
ha cambiado.
\end{observacion}

\subsubsection*{Combinación de ecuaciones}

Sustituyendo \(\theta_1 = \alpha_1 + \beta\) y \(\theta_2 = \alpha_2 + \beta\) en la ley
de Snell paraxial \(n_1\theta_1 = n_2\theta_2\):

\begin{align}
  n_1(\alpha_1 + \beta) &= n_2(\alpha_2 + \beta) \notag \\
  n_1\alpha_1 + n_1\beta &= n_2\alpha_2 + n_2\beta \notag \\
  n_1\alpha_1 - n_2\alpha_2 &= n_2\beta - n_1\beta \notag \\
  n_1\alpha_1 - n_2\alpha_2 &= (n_2 - n_1)\beta
  \label{eq:clave_virtual}
\end{align}

\subsubsection*{Expresiones geométricas y convención de signos}

En la aproximación de ángulos pequeños: \(\alpha_1 \approx h/p\), \(\alpha_2 \approx
h/|q|\), \(\beta \approx h/r\). Dado que \(q < 0\) para imagen virtual, entonces
\(|q| = -q\). Sustituyendo en \eqref{eq:clave_virtual} y dividiendo entre \(h\):

\[
  n_1 \cdot \frac{h}{p} - n_2 \cdot \frac{h}{|q|} = (n_2 - n_1) \cdot \frac{h}{r}
  \qquad \Longrightarrow \qquad
  \frac{n_1}{p} - \frac{n_2}{-q} = \frac{n_2 - n_1}{r}
\]

Es decir:

\begin{equation}
  \frac{n_1}{p} + \frac{n_2}{q} = \frac{n_2 - n_1}{r}
  \label{eq:fundamental_virtual}
\end{equation}

\textbf{Resultado:} la misma ecuación \eqref{eq:fundamental} para ambos casos. La
convención de signos (\(q > 0\) imagen real, \(q < 0\) imagen virtual) se encarga de
distinguirlos.

\subsection{Caso: Superficie cóncava con \texorpdfstring{\(n_1 < n_2\)}{n1 < n2}
            --- imagen virtual}

Un rayo \(\overrightarrow{PA}\) al refractarse en la superficie esférica cóncava
\textbf{diverge alejándose del eje}, pero si se prolonga hacia atrás intercepta el eje en
el punto \(Q\). De esta manera \(Q\) es la \textbf{imagen virtual} de \(P\).

La ley de Snell en aproximación paraxial es \(n_1 \cdot \theta_1 = n_2 \cdot \theta_2\).

De los triángulos \(\triangle APC\) y \(\triangle CAQ\) se obtiene (\(\beta\) es ángulo
exterior en ambos):

\[
  \beta = \alpha_1 + \theta_1 \;\Rightarrow\; \theta_1 = \beta - \alpha_1 \qquad;\qquad
  \beta = \alpha_2 + \theta_2 \;\Rightarrow\; \theta_2 = \beta - \alpha_2
\]

Sustituyendo en la ley de Snell paraxial \(n_1(\beta - \alpha_1) = n_2(\beta - \alpha_2)\)
y desarrollando:

\[
  n_1\alpha_1 - n_2\alpha_2 = -(n_2 - n_1)\beta
\]

Aplicando las aproximaciones de ángulos pequeños
(\(\alpha_1 \approx h/p\), \(\alpha_2 \approx h/q\), \(\beta \approx h/r\)) y dividiendo
entre \(h\):

\begin{equation}
  \frac{n_1}{p} - \frac{n_2}{q} = -\frac{(n_2 - n_1)}{r}
  \label{eq:concava_n1menor}
\end{equation}

\begin{observacion}
Nota: si aplicamos la convención de signos (\(q < 0\) y \(r < 0\) para superficie cóncava),
la ecuación vuelve a ser idéntica a la Ecuación Fundamental \eqref{eq:fundamental}:
\(\dfrac{n_1}{p} + \dfrac{n_2}{q} = \dfrac{n_2-n_1}{r}\).
\end{observacion}

\subsection{Caso: Superficie cóncava con \texorpdfstring{\(n_1 > n_2\)}{n1 > n2}
            --- imagen real}

Como la luz pasa a un medio menos denso (\(n_1 > n_2\)), el rayo se aleja de la normal.
Un rayo \(\overrightarrow{PA}\) al refractarse en la superficie esférica cóncava intercepta
al eje óptico directo en el punto \(Q\). De esta manera \(Q\) es la \textbf{imagen real}
de \(P\).

La ley de Snell en aproximación paraxial se mantiene: \(n_1 \cdot \theta_1 = n_2 \cdot
\theta_2\).

De los triángulos \(\triangle APC\) y \(\triangle CAQ\) (\(\beta\) es ángulo exterior):

\[
  \beta = \alpha_1 + \theta_1 \;\Rightarrow\; \theta_1 = \beta - \alpha_1 \qquad;\qquad
  \theta_2 = \alpha_2 + \beta \quad\text{(\(\theta_2\) es ángulo exterior en \(\triangle
  CAQ\))}
\]

Sustituyendo en la ley de Snell: \(n_1(\beta - \alpha_1) = n_2(\alpha_2 + \beta)\).
Desarrollando y reagrupando:

\[
  n_1\alpha_1 + n_2\alpha_2 = -(n_2 - n_1)\beta
\]

Aplicando \(\alpha_1 \approx h/p\), \(\alpha_2 \approx h/q\), \(\beta \approx h/r\) y
dividiendo entre \(h\):

\begin{equation}
  \frac{n_1}{p} + \frac{n_2}{q} = \frac{(n_2 - n_1)}{r}
  \label{eq:concava_n1mayor}
\end{equation}

\begin{observacion}
Conclusión final: si usas la convención de signos (\(r < 0\) para superficie cóncava),
\(q > 0\) por ser imagen real, llegas algebraicamente a la misma Ecuación Fundamental.
\end{observacion}

\textit{\textbf{Síntesis.}} Los cuatro casos de demostración —convexa/cóncava,
\(n_1 < n_2\)/\(n_1 > n_2\)— producen la misma ecuación algebraica cuando se aplica
rigurosamente la convención de signos. Esta universalidad no es trivial: es consecuencia
directa de la linearidad de la aproximación paraxial combinada con una convención
consistente. El resultado refleja la elegancia de la óptica gaussiana: una única expresión
\eqref{eq:fundamental} gobierna cuatro geometrías aparentemente distintas.

% ============================================================
\section{Convención de Signos}
% ============================================================

La convención de signos para los dioptrios supone que el \textbf{frente} de la superficie
es el lado desde el que se aproxima la luz.

\begin{itemize}
  \item \(p\) es \textbf{positivo} si el objeto está \textit{en frente} de la superficie
        (\textbf{objeto real}).
  \item \(p\) es \textbf{negativo} si el objeto está \textit{detrás} de la superficie
        (\textbf{objeto virtual}).
  \item \(q\) es \textbf{positivo} si la imagen está \textit{detrás} de la superficie
        (\textbf{imagen real}).
  \item \(q\) es \textbf{negativo} si la imagen está \textit{en frente} de la superficie
        (\textbf{imagen virtual}).
  \item \(r\) es \textbf{positivo} si el centro de curvatura está \textit{detrás} de la
        superficie (\textbf{superficie convexa}).
  \item \(r\) es \textbf{negativo} si el centro de curvatura está \textit{en frente} de la
        superficie (\textbf{superficie cóncava}).
\end{itemize}

Bajo esta estricta convención, la ecuación generalizada (Ec.\ 4-34) es:

\begin{equation}
  \frac{n_1}{p} + \frac{n_2}{q} = \frac{n_2 - n_1}{r}
  \label{eq:fundamental_general}
\end{equation}

\textit{\textbf{Síntesis.}} La convención de signos actúa como un ``codificador'' que
transforma la información geométrica (dónde está el objeto, dónde se forma la imagen, hacia
qué lado se curva la superficie) en información algebraica (el signo de \(p\), \(q\) y
\(r\)). Una vez interiorizada, elimina la necesidad de analizar cada caso geométrico por
separado.

% ============================================================
\section{Distancias Focales y Relación entre Focales}
% ============================================================

\subsection{Distancias focales}

Si el objeto está en el infinito (\(p \to \infty\)), la imagen se forma en el
\textbf{foco imagen} \(f'\):

\begin{equation}
  \frac{n_2}{f'} = \frac{n_2 - n_1}{r} \qquad\Longrightarrow\qquad
  f' = \frac{n_2 \cdot r}{n_2 - n_1}
  \label{eq:foco_imagen}
\end{equation}

Si la imagen está en el infinito (\(q \to \infty\)), el objeto está en el
\textbf{foco objeto} \(f\):

\begin{equation}
  \frac{n_1}{f} = \frac{n_2 - n_1}{r} \qquad\Longrightarrow\qquad
  f = \frac{n_1 \cdot r}{n_2 - n_1}
  \label{eq:foco_objeto}
\end{equation}

\subsection{Relación entre focales}

Dividiendo las expresiones de \(f'\) y \(f\):

\begin{equation}
  \frac{f'}{f} = \frac{\dfrac{n_2 \cdot r}{n_2 - n_1}}{\dfrac{n_1 \cdot r}{n_2 - n_1}}
  = \frac{n_2}{n_1}
  \label{eq:relacion_focales}
\end{equation}

\begin{observacion}
Obsérvese que las distancias focales \textbf{no son iguales}: \(f'/f = n_2/n_1\). Esto es
diferente a los espejos esféricos, donde ambas focales son iguales.
\end{observacion}

\textit{\textbf{Síntesis.}} La asimetría \(f' \neq f\) en los dioptrios refleja el hecho
de que la velocidad de la luz es distinta a ambos lados de la superficie. En un espejo, el
mismo medio está en ambos lados, de modo que \(f' = f\). En el dioptrio, los focos objeto
e imagen miden efectivamente la ``distancia óptica'' en medios con velocidades diferentes,
por lo que su cociente recoge exactamente la relación de índices.

% ============================================================
\section{Aumento Lateral de un Dioptrio}
% ============================================================

Para obtener el \textbf{aumento lateral} \(A\) de un dioptrio se utilizan dos rayos
principales a partir del objeto \(P\) de altura \(h\):

\begin{itemize}
  \item \textbf{Rayo por el centro (naranja):} se dirige hacia el centro de curvatura
        \(C\). Por ser normal a la superficie, no se desvía.
  \item \textbf{Rayo al vértice (azul):} incide exactamente en el vértice óptico \(O\).
        Allí la normal es el propio eje del sistema. Se refracta cruzando al otro rayo en
        \(Q\), formando la imagen de altura \(h'\).
\end{itemize}

\subsection{Geometría en el vértice}

Considerando únicamente el rayo refractado en el vértice \(O\), evaluamos los dos
triángulos rectángulos correspondientes formados con el eje (\(\triangle PAO\) y
\(\triangle BOQ\)):

\[
  \tan\alpha_1 = \frac{h}{p} \qquad;\qquad \tan\alpha_2 = -\frac{h'}{q}
\]

(El signo negativo en \(\alpha_2\) refleja que la imagen está invertida geométricamente.)

De acuerdo con la ley de Snell aplicada en el punto \(O\) (donde el eje es la normal
recta):
\[
  n_1 \cdot \sin\alpha_1 = n_2 \cdot \sin\alpha_2
\]

Si los ángulos de incidencia son suficientemente pequeños (rayos cercanos al eje):
\[
  \tan\alpha_1 \approx \sin\alpha_1 \qquad\text{y}\qquad \tan\alpha_2 \approx \sin\alpha_2
\]

Por lo tanto:
\[
  n_1 \tan\alpha_1 = n_2 \tan\alpha_2 \qquad\Longrightarrow\qquad
  n_1 \left(\frac{h}{p}\right) = n_2 \left(-\frac{h'}{q}\right)
\]

\subsection{Expresión final del aumento}

El aumento lateral se define por la proporción de alturas \(A = h'/h\). Despejando:

\begin{equation}
  \boxed{A = \frac{h'}{h} = -\frac{n_1 q}{n_2 p}}
  \label{eq:aumento}
\end{equation}

\begin{observacion}
El signo negativo inherente indica que si la imagen es real (\(q > 0\)), el aumento resulta
negativo y la imagen estará invertida respecto al objeto.
\end{observacion}

\textit{\textbf{Síntesis.}} La fórmula del aumento lateral \eqref{eq:aumento} codifica
simultáneamente tres informaciones: el módulo (\(|A|\)) indica cuánto se amplifica o
reduce la imagen; el signo indica si está invertida (\(A < 0\)) o derecha (\(A > 0\)); y
la presencia de \(n_1/n_2\) muestra que el aumento depende no solo de la geometría
(\(q/p\)), sino también de los medios ópticos.

% ============================================================
\section{Superficies Refractoras Planas}
% ============================================================

Si la superficie refractora es plana, el radio de curvatura es infinito matemáticamente
(\(r \to \infty\)). Bajo esta condición, la Ecuación Fundamental \eqref{eq:fundamental_general}
se reduce drásticamente porque su lado derecho se aproxima a cero:

\[
  \frac{n_1}{p} + \frac{n_2}{q} = \frac{n_2 - n_1}{\infty} = 0
  \qquad\Longrightarrow\qquad \frac{n_1}{p} = -\frac{n_2}{q}
\]

O bien, resolviendo de forma explícita para la distancia de la imagen \(q\):

\begin{equation}
  \boxed{q = -\left(\frac{n_2}{n_1}\right)p}
  \label{eq:dioptrio_plano}
\end{equation}

\subsubsection*{Análisis de posición e imagen}

De la ecuación \eqref{eq:dioptrio_plano} vemos que el signo de \(q\) siempre es opuesto
al de \(p\).

\begin{observacion}
La imagen formada por una superficie refractante plana está siempre en el mismo lado de la
superficie que el objeto original. Constituye siempre una \textbf{imagen virtual}.
\end{observacion}

En el caso típico en que \(n_1 > n_2\) (como mirar agua desde el aire), donde se forma una
imagen virtual entre el objeto y la superficie. Si \(n_1 < n_2\), la imagen sigue siendo
virtual, pero se forma a la izquierda (más profunda) del objeto original.

\begin{observacion}
Dato interesante: si combinas la fórmula del aumento lateral \eqref{eq:aumento} con el
valor de \(q\) obtenido aquí \eqref{eq:dioptrio_plano}, se cancelan las variables y se
descubre que \(A = +1\). Es decir, la imagen en un dioptrio plano se ve del mismo tamaño
y nunca invertida.
\end{observacion}

\textit{\textbf{Síntesis.}} El dioptrio plano es, paradójicamente, el caso más cotidiano
de dioptrio: toda interfaz plana entre medios ópticos distintos —la superficie de una
piscina, el vidrio de un acuario, la córnea en contacto con el aire— actúa como dioptrio
plano. Su consecuencia más llamativa es la ilusión de profundidad reducida: un objeto a
profundidad real \(h\) parece estar a profundidad aparente \((n_2/n_1)\,h\), siempre menor
cuando se mira desde el medio menos denso.

% ============================================================
\section{Problemas Resueltos}
% ============================================================

\begin{ejemplo}[Varilla de vidrio en el aire (Ejemplo 4-18)]
Una varilla de vidrio cilíndrica en el aire tiene un índice de refracción de \(n = 1{,}52\).
Se puló un extremo para formar una superficie semiesférica con radio \(r = 2{,}00\) cm.

\begin{itemize}
  \item[\textbf{a)}] Calcule la distancia de la imagen de un objeto pequeño situado sobre
        el eje de la varilla, a 8,00 cm a la izquierda del vértice.
  \item[\textbf{b)}] Obtenga el aumento lateral.
\end{itemize}

\textit{Solución (Parte a):}

Identificamos los valores bajo la convención de signos:
\begin{itemize}
  \item \(n_1 = 1{,}00\) (el objeto está en el aire)
  \item \(n_2 = 1{,}52\) (la luz ingresa al vidrio)
  \item \(p = +8{,}00\) cm (objeto en frente de la superficie)
  \item \(r = +2{,}00\) cm (superficie convexa, centro detrás del vértice)
\end{itemize}

Aplicando la Ecuación Fundamental \eqref{eq:fundamental_general}:

\[
  \frac{1}{8\;\text{cm}} + \frac{1{,}52}{q} = \frac{1{,}52 - 1}{2\;\text{cm}}
\]

Despejando \(q\):

\begin{align*}
  \frac{1{,}52}{q} &= \frac{0{,}52}{2} - \frac{1}{8} = 0{,}26 - 0{,}125 = 0{,}135\;\text{cm}^{-1} \\
  q &= +11{,}3\;\text{cm} \quad \textbf{(Imagen real)}
\end{align*}

\textit{Solución (Parte b):}

Calculamos el aumento lateral con la fórmula \eqref{eq:aumento}:

\[
  A = -\frac{n_1 q}{n_2 p} = -\frac{1}{1{,}52}\left(\frac{11{,}3\;\text{cm}}{8\;\text{cm}}\right)
  = -0{,}929
\]

El valor \(A = -0{,}929\) indica que la imagen es un poco más pequeña que el objeto
original y que la imagen está \textbf{invertida}. Por ejemplo, si el objeto fuera una
flechita de 1,00 mm apuntando hacia arriba, la imagen sería una flechita de aprox.\ 0,93
mm apuntando hacia abajo.
\end{ejemplo}

\begin{ejemplo}[Varilla de vidrio sumergida en agua (Ejemplo 4-19)]
Se sumerge en agua (índice de refracción \(n_1 = 1{,}33\)) la varilla de vidrio del
ejemplo anterior tal como se muestra. Las demás magnitudes tienen los mismos valores que
en el caso anterior. Obtenga la distancia de imagen y el aumento lateral.

\textit{Solución (Parte a):}

\begin{itemize}
  \item \(n_1 = 1{,}33\) (el objeto ahora está en el agua)
  \item \(n_2 = 1{,}52\) (la luz ingresa al vidrio)
  \item \(p = +8{,}00\) cm (distancia del objeto al vértice convexo)
  \item \(r = +2{,}00\) cm (se mantiene la geometría del radio)
\end{itemize}

\[
  \frac{1{,}33}{8\;\text{cm}} + \frac{1{,}52}{q} = \frac{(1{,}52 - 1{,}33)}{+2\;\text{cm}}
\]

\[
  0{,}16625 + \frac{1{,}52}{q} = 0{,}095
\]

\[
  \frac{1{,}52}{q} = 0{,}095 - 0{,}16625 = -0{,}07125\;\text{cm}^{-1}
\]

\[
  q = -21{,}3\;\text{cm} \quad \textbf{(Imagen virtual)}
\]

\textit{Solución (Parte b):}

\[
  A = -\frac{1{,}33}{1{,}52}\left(\frac{-21{,}3\;\text{cm}}{8\;\text{cm}}\right) = +2{,}32
\]

Al cambiar de aire a agua, hemos reducido drásticamente la potencia refractiva de la
superficie porque la diferencia de índices \((1{,}52 - 1{,}33)\) es mucho menor. Como
resultado, la superficie ya no logró enfocar los rayos y formó una imagen virtual, la cual
no está invertida y ahora se ve más del doble de grande \((2{,}32\times)\).
\end{ejemplo}

\begin{ejemplo}[Superficie cóncava (Ejemplo 4-20)]
El extremo izquierdo de una larga varilla de vidrio (\(n = 1{,}50\)) se esmerila para
formar una \textbf{superficie hemisférica cóncava} de radio \(r = 4{,}00\) cm. Un objeto
en forma de flecha de 3,0 mm de altura se sitúa sobre el eje a 16,0 cm del vértice.
Calcule la posición, la altura de la imagen, y determine si es derecha o invertida.

\textit{Solución (Parte a: Posición):}

\begin{itemize}
  \item \(n_1 = 1{,}00\) (Aire), \(n_2 = 1{,}50\) (Vidrio)
  \item \(p = +16{,}0\) cm (objeto en frente del dioptrio)
  \item \(r = -4{,}00\) cm (superficie CÓNCAVA: centro de curvatura en el frente, \(r\)
        negativo)
\end{itemize}

\[
  \frac{1}{16\;\text{cm}} + \frac{1{,}5}{q} = \frac{(1{,}5 - 1)}{-4\;\text{cm}}
\]

\[
  \frac{1{,}5}{q} = -0{,}125 - 0{,}0625 = -0{,}1875\;\text{cm}^{-1}
\]

\[
  q = -8{,}00\;\text{cm} \quad \textbf{(Imagen virtual)}
\]

\textit{Solución (Parte b: Aumento):}

\[
  A = -\frac{1}{1{,}5}\left(\frac{-8{,}00\;\text{cm}}{16{,}0\;\text{cm}}\right)
  = -\frac{1}{1{,}5}(-0{,}5) = +\frac{1}{3} \approx +0{,}333
\]

Como \(A\) es positivo, la imagen es \textbf{derecha}. Su altura final será:
\[
  h' = A \cdot h = \left(\frac{1}{3}\right)(3{,}0\;\text{mm}) = 1{,}0\;\text{mm}
\]

El signo negativo en \(q\) confirma que la imagen es virtual (se forma a la izquierda, en
el medio de índice menor). El signo positivo en \(A\) demuestra matemáticamente que la
imagen es derecha (apuntando hacia arriba, igual que el objeto) pero reducida exactamente
a un tercio de la original.
\end{ejemplo}

\begin{ejemplo}[El ojo del pez --- Pecera (Ejemplo 4-21)]
Un pez dentro de una pecera esférica observa un gato situado en el exterior. La pecera
tiene un diámetro de 80,0 cm (\(r = 40\) cm). Los índices del aire y agua son 1,00 y
1,33. La distancia del vértice de la pecera al gato es \(p = 50{,}0\) cm.

Determina: \textbf{a)} la posición de la imagen del gato respecto al vértice de la pecera;
\textbf{b)} el aumento lateral de la imagen sugerido matemáticamente.

\textit{Solución (Parte a: Posición):}

Aplicando la Convención de Signos (la luz viaja desde el gato hacia el pez):
\begin{itemize}
  \item \(n_1 = 1{,}00\) (Aire, desde donde viene la luz del objeto gato)
  \item \(n_2 = 1{,}33\) (Agua, hacia donde va la luz)
  \item \(p = +50{,}0\) cm (el objeto gato está en frente del vértice)
  \item \(r = +40{,}0\) cm (al ser convexa para la luz entrante del gato, el centro de
        curvatura \(C\) se encuentra detrás del vértice)
\end{itemize}

\[
  \frac{1}{50\;\text{cm}} + \frac{1{,}33}{q} = \frac{(1{,}33 - 1)}{+40\;\text{cm}}
\]

\[
  0{,}020 + \frac{1{,}33}{q} = 0{,}00825
\]

\[
  \frac{1{,}33}{q} = 0{,}00825 - 0{,}020 = -0{,}01175\;\text{cm}^{-1}
\]

\[
  q \approx -113{,}2\;\text{cm} \quad \textbf{(Imagen virtual)}
\]

\textit{Solución (Parte b: Aumento):}

Usamos la Ecuación \eqref{eq:aumento} evaluando con el \(q\) negativo recién obtenido:

\[
  A = -\frac{n_1 q}{n_2 p} = -\frac{1}{1{,}33}\left(\frac{-113{,}2\;\text{cm}}{50{,}0\;\text{cm}}\right)
  = +1{,}70
\]

Como \(q\) es negativo, la imagen del gato se forma a 113,2 cm a la izquierda del vidrio
(mucho más atrás de donde está el gato real en 50 cm). El aumento lateral es \(+1{,}70\),
es decir, la imagen es \textbf{derecha} y un 70\% más grande. Para el pobre pececito,
¡el gato se ve como un enorme monstruo de aspecto virtual!
\end{ejemplo}

\begin{ejemplo}[Profundidad aparente --- La moneda en la piscina (Ejemplo clásico)]
Analizamos la profundidad aparente de una piscina que contiene agua (índice
\(n_1 = 1{,}33 \approx 4/3\)) de altura real \(h\), si se mira directamente desde arriba
desde el aire (\(n_2 = 1{,}00\)).

Como la superficie del agua es perfectamente \textbf{plana}, podemos utilizar directamente
la simplificación del dioptrio plano \eqref{eq:dioptrio_plano} evaluada con los valores
del problema donde la distancia del objeto \(P\) es la profundidad geométrica real de la
piscina (\(p = h\)):

\begin{itemize}
  \item \(n_1 = 4/3\) (Agua, el medio donde origina la luz la moneda)
  \item \(n_2 = 1\) (Aire, hacia donde llega la luz a nuestros ojos)
  \item \(p = h\) (Profundidad real)
\end{itemize}

\[
  q = -\left(\frac{n_2}{n_1}\right)p = -\left(\frac{1}{4/3}\right)h = -\left(\frac{3}{4}\right)h
\]

\begin{equation}
  \boxed{q = -\frac{3}{4}\,h}
  \label{eq:profundidad_aparente}
\end{equation}

\textbf{Conclusión de diseño:} la profundidad aparente calculada \(\left(-\frac{3}{4}h\right)\)
es solamente las tres cuartas partes de la profundidad real de la piscina.

La distancia de la imagen resultó negativa. De acuerdo con las reglas de signos, esto
significa que la imagen es virtual y se forma obligatoriamente en el lado entrante de la
superficie plana (es decir, sumergida en el mismo medio del objeto), dándonos la famosa
ilusión óptica de la poca profundidad.
\end{ejemplo}

\textit{\textbf{Síntesis.}} Los cinco ejemplos resueltos ilustran la versatilidad de la
ecuación fundamental \eqref{eq:fundamental_general}: un mismo formalismo algebraico
describe desde la microscopía (varilla de vidrio) hasta la percepción visual en la vida
cotidiana (piscina, acuario). La clave pedagógica es que el trabajo se concentra en la
correcta asignación de signos; la aritmética es secundaria. El signo de \(q\) al final
actúa como veredicto físico: real o virtual, dentro o fuera, invertida o derecha.

% ============================================================
\section*{Ecuaciones Clave}
% ============================================================

\begin{center}
\begin{tabular}{@{}clll@{}}
\toprule
\textbf{\#} & \textbf{Descripción} & \textbf{Ecuación} & \textbf{Ref.} \\
\midrule
1 & Ley de Snell (paraxial) &
  \(\displaystyle n_1\theta_1 = n_2\theta_2\) &
  \eqref{eq:snell_paraxial} \\[6pt]
2 & Ecuación fundamental del dioptrio esférico &
  \(\displaystyle \dfrac{n_1}{p} + \dfrac{n_2}{q} = \dfrac{n_2 - n_1}{r}\) &
  \eqref{eq:fundamental_general} \\[10pt]
3 & Foco imagen &
  \(\displaystyle f' = \dfrac{n_2\, r}{n_2 - n_1}\) &
  \eqref{eq:foco_imagen} \\[10pt]
4 & Foco objeto &
  \(\displaystyle f = \dfrac{n_1\, r}{n_2 - n_1}\) &
  \eqref{eq:foco_objeto} \\[10pt]
5 & Relación entre focales &
  \(\displaystyle \dfrac{f'}{f} = \dfrac{n_2}{n_1}\) &
  \eqref{eq:relacion_focales} \\[10pt]
6 & Aumento lateral &
  \(\displaystyle A = -\dfrac{n_1\, q}{n_2\, p}\) &
  \eqref{eq:aumento} \\[10pt]
7 & Dioptrio plano &
  \(\displaystyle q = -\dfrac{n_2}{n_1}\,p\) &
  \eqref{eq:dioptrio_plano} \\[10pt]
8 & Profundidad aparente &
  \(\displaystyle q_{\text{ap}} = -\dfrac{n_2}{n_1}\,h\) &
  \eqref{eq:profundidad_aparente} \\
\bottomrule
\end{tabular}
\end{center}

\end{document}
